\section{The Blaze PCS}

We can now assemble the pieces.  The PCS witness is a multilinear polynomial
represented by its Boolean evaluation tensor
\[
  m\in\F^{\B^\ell\times\B^k}.
\]
The first coordinate \(h\in\B^\ell\) selects one slice, and the second
coordinate \(a\in\B^k\) selects a coordinate inside that slice.  The verifier's
opening claim is
\[
  z=(z_{\mathrm{slice}},z_{\mathrm{in}})\in\F^\ell\times\F^k,
  \qquad
  v\in\F,
  \qquad
  \MLE(m)(z)=v.
\]

Blaze does not authenticate this whole tensor by running a foldable-code PCS on
all \(2^{\ell+k}\) entries.  It first encodes each slice with the fast RAA code
\[
  \RAA:\F^{\B^k}\to\F^{\B^n},
  \qquad K=2^k,\quad N=2^n,\quad k<n,
\]
and commits to the encoded matrix
\[
  c\in\F^{\B^\ell\times\B^n},
  \qquad
  c[h,b]=\RAA(m[h,\ast])[b].
\]
This gives the verifier coordinate access to encoded columns \(c[\ast,b]\).
The row-fold step then compresses all slices into one virtual codeword.  The
key word is virtual: after the fold, the verifier has an openable tensor
\(c_\star\), but \(c_\star\) is not promised to be an RAA codeword.  The later
native RAA checks and \(\MLEAuth\) backend are what authenticate the claim that
the folded object is close to an honest RAA encoding of a folded message.

\subsection{The Row-Fold Virtual Handle}

The row-fold handle implements the message-oracle interface from
Figure~\ref{fig:func-msg-oracle}.  It does not store a new vector.  It stores an
existing handle for the committed encoded matrix and the verifier's fold vector
\(\rho\in\F^{\B^\ell}\).  When a backend asks for \(c_\star[b]\), the handle
opens the already-committed column \(c[\ast,b]\) and returns its inner product
with \(\rho\).

% DO NOT MODIFY: Frozen semantic anchor for the Blaze PCS.
% This figure defines the only row-fold commitment topology used by Figure~\ref{fig:protocol-blaze-pcs}.
\begin{protocolbox}{Virtual Handle: Row Fold}{fig:protocol-row-fold-handle}
\scriptsize\raggedright
\noindent\textbf{DO NOT MODIFY.} Frozen semantic anchor for the Blaze PCS.

\smallskip
\noindent Public parameters: slice domain \(\B^\ell\), codeword domain
\(\B^n\), and the message-oracle functionality
\(\mathcal F_{\mathrm{MsgOracle}}\).

\smallskip
\noindent\(\underline{\mathsf{RowFold.Commit}(\verifier:(R_c,\rho\in\F^{\B^\ell}),\ \prover:(R_c,\rho))\to R_\star\in\Handle}\)

\smallskip
\noindent Here \(R_c\) is a message-oracle handle for a tensor
\(C\in(\F^{\B^\ell})^{\B^n}\).

\smallskip
\begin{enumerate}[leftmargin=0.24in,label=\textbf{F\arabic*.},itemsep=0.06em,topsep=0.06em,parsep=0pt]
  \item \(\mathsf{RowFold}\): create a fresh handle \(R_\star\).
  \item \(\mathsf{RowFold}\): store \((R_c,\rho)\) under \(R_\star\).
  \item \(\mathsf{RowFold}\): return \(R_\star\).
\end{enumerate}

\smallskip
\noindent\(\underline{\mathsf{RowFold.Open}(\verifier:(R_\star,b\in\B^n))\to(\verifier:o\in\F\cup\{\bot\})}\)

\smallskip
\begin{enumerate}[leftmargin=0.24in,label=\textbf{G\arabic*.},itemsep=0.06em,topsep=0.06em,parsep=0pt]
  \item \(\mathsf{RowFold}\): read \((R_c,\rho)\) stored under \(R_\star\).
  \item \(\mathsf{RowFold}\): call
        \(\mathcal F_{\mathrm{MsgOracle}}.\mathsf{Open}(\verifier:(R_c,b))\),
        receiving \(w\in\F^{\B^\ell}\cup\{\bot\}\).
  \item \(\mathsf{RowFold}\): if \(w=\bot\), return \(\bot\).
  \item \(\mathsf{RowFold}\): return
        \((\verifier:o=\sum_{h\in\B^\ell}\rho[h]w[h])\).
\end{enumerate}
\end{protocolbox}

The tensor opened by \(R_\star\) is therefore
\[
  c_\star[b]=\sum_{h\in\B^\ell}\rho[h]c[h,b],
  \qquad b\in\B^n.
\]
This definition is binding because \(R_c\) was already fixed before
\(\rho\) was sampled.

\subsection{Compiled RAA Adapter}

The native RAA flow in Figure~\ref{fig:native-blaze-fmliop-hybrid} is written in the
MLIOP model, where a native oracle handle \(O_H\) is a source of MLE values.  A
real compiled protocol has coordinate-opening handles \(C_H\) instead.  The
adapter below is the bridge: it follows the same algebra and timing as the
native RAA compiler, but every native evaluation request is later authenticated
through \(\MLEAuth\).

Equivalently, the adapter turns native requests into RMLE-style authentication
tasks.  For one tensor \(H\in\F^{\B^d}\), an RMLE task consists of a
coordinate-opening handle \(C_H\) for the Boolean entries of \(H\) together
with an MLE claim
\[
  \MLE(H)(r)=v,\qquad r\in\F^d.
\]
The backend may encode \(H\), prove proximity to that encoded object, and
answer the MLE claim through BaseFold/RMLE.  The RAA algebra does not inspect
that backend encoding; it only supplies the list of claims to authenticate.

We first name the compiled versions of the two reducers.  These are not new
algebraic checks.  They are the same accumulator and permutation reductions
from Figures~\ref{fig:native-accumulate-reduce} and
\ref{fig:native-permute-reduce}, but the handle component in every request is a
coordinate-opening handle \(C_H\), not a native MLIOP handle \(O_H\).  The
resulting \(\EvalReqs\) are therefore ready to be translated into \(\MLEAuth\)
lanes.
The compiled \(u_1\) coordinate handle is virtual in the same sense as
Figure~\ref{fig:def-virtual-oracle-opening}: the box gives its opening rule
directly.

% DO NOT MODIFY: Frozen semantic anchor for the Blaze PCS.
% These reducers are helper figures for Figure~\ref{fig:protocol-compiled-raa-check}.
\begin{protocolbox}{Subroutine: Compiled RAA Reducers}{fig:protocol-compiled-raa-reducers}
\scriptsize\raggedright
\noindent\textbf{DO NOT MODIFY.} Frozen semantic anchor for the Blaze PCS.

\smallskip
\noindent Public parameters: \(n\), \(\widehat A:\F^n\times\F^n\to\F\),
\(\Idx,\widehat\pi:\F^n\to\F\), and product-tree maps
\(\mathsf{root},\mathsf{leaf},\mathsf{parent},\mathsf{left},\mathsf{right}\).

\smallskip
\noindent\(\underline{\mathsf{CompiledAccumulate.Reduce}(\verifier:(H_x,C_x,H_y,C_y,r\in\F^n),\ \prover:(x,y\in\F^{\B^n}))\to(\EvalReqs,\SumClaims)}\)

\smallskip
\begin{enumerate}[leftmargin=0.24in,label=\textbf{CA\arabic*.},itemsep=0.06em,topsep=0.06em,parsep=0pt]
  \item \(\prover\to\verifier\): \(s\in\F\).
  \item \(\verifier\): \(\EvalReqs\gets\{(H_y,C_y,r,s)\}\).
  \item \(\verifier\): \(\SumClaims\gets\bigl((s,H_x,C_x,G(X)=\widehat A(r,X)\widehat H_x(X))\bigr)\).
  \item \(\verifier\): return \((\EvalReqs,\SumClaims)\).
\end{enumerate}

\smallskip
\noindent\(\underline{\mathsf{CompiledPermute.Reduce}(\verifier:(H_x,C_x,H_y,C_y,H_f,C_f,H_g,C_g,\widehat\pi,\alpha,\beta,r\in\F^n),\ \prover:(x,y,f,g))\to(\EvalReqs,\mathcal T,\SumClaims)}\)

\smallskip
\begin{enumerate}[leftmargin=0.24in,label=\textbf{CP\arabic*.},itemsep=0.06em,topsep=0.06em,parsep=0pt]
  \item \(\prover\to\verifier\): \(v_{f,\mathrm{root}},v_{g,\mathrm{root}},
        v_x,v_y,v_f,v_g\in\F\).
  \item \(\verifier\): \(\EvalReqs\gets
        \{(H_f,C_f,\mathsf{root},v_{f,\mathrm{root}}),
        (H_g,C_g,\mathsf{root},v_{g,\mathrm{root}})\}\).
  \item \(\verifier\): \(\EvalReqs\gets\EvalReqs\cup
        \{(H_x,C_x,r,v_x),(H_y,C_y,r,v_y),
        (H_f,C_f,\mathsf{leaf}(r),v_f),(H_g,C_g,\mathsf{leaf}(r),v_g)\}\).
  \item \(\verifier\): \(\mathcal T\gets
        \{v_{f,\mathrm{root}}=v_{g,\mathrm{root}},
        v_f=\alpha-(v_x+\beta\,\widehat\pi(r)),
        v_g=\alpha-(v_y+\beta\,\Idx(r))\}\).
  \item \(\verifier\): \(\SumClaims\gets\bigl((0,H_f,C_f,G_f(X)=\chi_r(X)\mathsf{TreeRes}(H_f,X)),(0,H_g,C_g,G_g(X)=\chi_r(X)\mathsf{TreeRes}(H_g,X))\bigr)\).
  \item \(\verifier\): return \((\EvalReqs,\mathcal T,\SumClaims)\).
\end{enumerate}
\end{protocolbox}

% DO NOT MODIFY: Frozen semantic anchor for the Blaze PCS.
% This figure is the compiled RAA adapter called by Figure~\ref{fig:protocol-blaze-pcs}.
\begin{protocolbox}{Subroutine: Compiled RAA Check}{fig:protocol-compiled-raa-check}
\scriptsize\raggedright
\noindent\textbf{DO NOT MODIFY.} Frozen semantic anchor for the Blaze PCS.

\smallskip
\noindent Public parameters: \(k,n\), the RAA stage maps, the compiled reducers
from Figure~\ref{fig:protocol-compiled-raa-reducers}, the handle-parametric six-sumcheck
from Figure~\ref{fig:native-six-sumcheck}, the \(\MLEAuth\) implementation
from Figure~\ref{fig:protocol-mle-auth-basefold},
\(\mathcal F_{\mathrm{MsgOracle}}\), the repetition projection
\(\mathsf{msg}:\F^n\to\F^k\), \(\mathcal H_{\RAA}=\{m_\star,u_1,u_2,u_3,u_4,c_\star,f_1,g_1,f_2,g_2\}\), injective \(q_{\RAA}:\mathcal H_{\RAA}\to\mathcal L_{\RAA}\subseteq\B^4\), and \(\MLEAuth\) dimension \(D=n+1\).

\smallskip
\noindent\(\underline{\mathsf{CompiledRAA.Check}(\verifier:(C_{c_\star},z_{\mathrm{in}}\in\F^k,v_\star\in\F),\ \prover:(C_{c_\star},m_\star\in\F^{\B^k},c_\star^{\mathrm{wit}}\in\F^{\B^n}))\to(\verifier:a\in\{0,1\})}\)

\smallskip
\begin{enumerate}[leftmargin=0.24in,label=\textbf{A\arabic*.},itemsep=0.06em,topsep=0.06em,parsep=0pt]
  \item \(\prover\): choose \(u_2,u_3,u_4\in\F^{\B^n}\).
  \item \(\prover\): set \(h_{m_\star}=m_\star\), \(h_{u_i}=u_i\) for
        \(i\in\{2,3,4\}\), and \(h_{c_\star}=c_\star^{\mathrm{wit}}\).
  \item \(\prover,\verifier\): for
        \(H\in\{m_\star,u_2,u_3,u_4\}\), set
        \(C_H\gets\mathcal F_{\mathrm{MsgOracle}}.\mathsf{Submit}
        (\prover:h_H,\verifier:\bot)\).
  \item \(\prover,\verifier\): set \(C_{c_\star}\) to the handle supplied in
        the verifier input.
  \item \(\verifier\): sample \(\alpha,\beta\leftarrow\F\), then send
        \((\alpha,\beta)\) to \(\prover\).
  \item \(\prover\): define the repetition tensor
        \(u_1=F_{\mathrm{rep}}(m_\star)\in\F^{\B^n}\).
  \item \(\prover\): set \(h_{u_1}=u_1\).
  \item \(\prover,\verifier\): create virtual message-oracle handle \(C_{u_1}\).
  \item \(\prover,\verifier\): set \(\mathsf{Open}(C_{u_1},b):=\mathsf{Open}(C_{m_\star},\mathsf{msg}(b))\) for \(b\in\B^n\).
  \item \(\prover\): construct \(f_1,g_1\in\F^{\B^{n+1}}\) by
        Steps P1--P5 of Figure~\ref{fig:native-permute-reduce}, with
        \((x,y,\widehat\pi)=(u_1,u_2,\widehat\pi_1)\).
  \item \(\prover\): construct \(f_2,g_2\in\F^{\B^{n+1}}\) by
        Steps P1--P5 of Figure~\ref{fig:native-permute-reduce}, with
        \((x,y,\widehat\pi)=(u_3,u_4,\widehat\pi_2)\).
  \item \(\prover\): set \(h_T=T\) for
        \(T\in\{f_1,g_1,f_2,g_2\}\).
  \item \(\prover,\verifier\): for
        \(T\in\{f_1,g_1,f_2,g_2\}\), set
        \(C_T\gets\mathcal F_{\mathrm{MsgOracle}}.\mathsf{Submit}
        (\prover:T,\verifier:\bot)\).
  \item \(\prover,\verifier\): for each \(H\in\mathcal H_{\RAA}\), define
        \(h_{q_{\RAA}(H)}=h_H\), \(C_{q_{\RAA}(H)}=C_H\), with dimensions
        \(d_{q_{\RAA}(m_\star)}=k\),
        \(d_{q_{\RAA}(u_i)}=d_{q_{\RAA}(c_\star)}=n\) for
        \(i\in\{1,2,3,4\}\), and
        \(d_{q_{\RAA}(f_i)}=d_{q_{\RAA}(g_i)}=n+1\) for \(i\in\{1,2\}\).
  \item \(\prover,\verifier\): \(A_{\RAA}\gets
        \MLEAuth.\mathsf{Commit}
        (\prover:(h_q)_{q\in\mathcal L_{\RAA}},
        \verifier:(d_q,C_q)_{q\in\mathcal L_{\RAA}})\).
  \item \(\verifier\): sample \(r_{\mathrm{chk}}\leftarrow\F^n\), then send
        \(r_{\mathrm{chk}}\) to \(\prover\).
  \item \(\prover\to\verifier\): \(v_m\in\F\).
  \item \(\verifier\): \(\EvalReqs_{\mathrm{val}}\gets
        \{(m_\star,C_{m_\star},z_{\mathrm{in}},v_m)\}\) and
        \(\mathcal T_{\mathrm{val}}\gets\{v_m=v_\star\}\).
  \item \(\prover,\verifier\):
        \((\EvalReqs_{\mathrm{acc},1},\SumClaims_{\mathrm{acc},1})\gets
        \mathsf{CompiledAccumulate.Reduce}
        (\verifier:(H_{u_2},C_{u_2},H_{u_3},C_{u_3},r_{\mathrm{chk}}),
        \prover:(u_2,u_3))\).
  \item \(\prover,\verifier\):
        \((\EvalReqs_{\mathrm{acc},2},\SumClaims_{\mathrm{acc},2})\gets
        \mathsf{CompiledAccumulate.Reduce}
        (\verifier:(H_{u_4},C_{u_4},H_{c_\star},C_{c_\star},r_{\mathrm{chk}}),
        \prover:(u_4,c_\star^{\mathrm{wit}}))\).
  \item \(\prover,\verifier\):
        \((\EvalReqs_{\mathrm{perm},1},\mathcal T_{\mathrm{perm},1},
        \SumClaims_{\mathrm{perm},1})\gets
        \mathsf{CompiledPermute.Reduce}
        (\verifier:(H_{u_1},C_{u_1},H_{u_2},C_{u_2},H_{f_1},C_{f_1},H_{g_1},C_{g_1},
        \widehat\pi_1,\alpha,\beta,r_{\mathrm{chk}}),
        \prover:(u_1,u_2,f_1,g_1))\).
  \item \(\prover,\verifier\):
        \((\EvalReqs_{\mathrm{perm},2},\mathcal T_{\mathrm{perm},2},
        \SumClaims_{\mathrm{perm},2})\gets
        \mathsf{CompiledPermute.Reduce}
        (\verifier:(H_{u_3},C_{u_3},H_{u_4},C_{u_4},H_{f_2},C_{f_2},H_{g_2},C_{g_2},
        \widehat\pi_2,\alpha,\beta,r_{\mathrm{chk}}),
        \prover:(u_3,u_4,f_2,g_2))\).
  \item \(\verifier\): \(\EvalReqs_{\mathrm{red}}\gets
        \EvalReqs_{\mathrm{acc},1}\cup\EvalReqs_{\mathrm{acc},2}\cup
        \EvalReqs_{\mathrm{perm},1}\cup\EvalReqs_{\mathrm{perm},2}\).
  \item \(\verifier\): let \(B_1\gets\SumClaims_{\mathrm{acc},1}[1]\), \(B_2\gets\SumClaims_{\mathrm{acc},2}[1]\), \(B_3\gets\SumClaims_{\mathrm{perm},1}[1]\), \(B_4\gets\SumClaims_{\mathrm{perm},1}[2]\), \(B_5\gets\SumClaims_{\mathrm{perm},2}[1]\), \(B_6\gets\SumClaims_{\mathrm{perm},2}[2]\).
  \item \(\verifier\): \(\SumClaims_{\RAA}\gets(B_1,B_2,B_3,B_4,B_5,B_6)\).
  \item \(\verifier\): \(\mathcal T_{\mathrm{red}}\gets
        \mathcal T_{\mathrm{perm},1}\cup\mathcal T_{\mathrm{perm},2}\).
  \item \(\prover,\verifier\): run Figure~\ref{fig:native-six-sumcheck} on
        \((\SumClaims_{\RAA},r_{\mathrm{chk}})\), obtaining
        \((a_{\mathrm{sc}},\EvalReqs_{\mathrm{sc}})\).
  \item \(\verifier\): if \(a_{\mathrm{sc}}=0\), return \((\verifier:0)\).
  \item \(\verifier\): if any equality in
        \(\mathcal T_{\mathrm{val}}\cup\mathcal T_{\mathrm{red}}\) is false,
        return \((\verifier:0)\).
  \item \(\verifier\): \(\EvalReqs_{\mathrm{auth}}\gets
        \EvalReqs_{\mathrm{val}}\cup\EvalReqs_{\mathrm{red}}\cup
        \EvalReqs_{\mathrm{sc}}\).
  \item \(\verifier\): translate every
        \((H,C_H,r,v_H)\in\EvalReqs_{\mathrm{auth}}\) into
        \((q_{\RAA}(H),r,v_H)\in\mathcal E_{\RAA}\).
  \item \(\prover,\verifier\): \(a_{\mathrm{auth}}\gets
        \MLEAuth.\mathsf{Eval}
        (\verifier:(A_{\RAA},\mathcal E_{\RAA}),
        \prover:(h_q)_{q\in\mathcal L_{\RAA}})\).
  \item \(\verifier\): return \((\verifier:a_{\mathrm{auth}})\).
\end{enumerate}
\end{protocolbox}

Before we discuss batching these claims, it is useful to spell out exactly what
the compiled RAA check has produced.  Every authentication request has the form
\[
  (H,C_H,r,v_H),
  \qquad
  H\in\mathcal H_{\RAA},\quad r\in\F^{d_H},\quad
  \MLE(H)(r)=v_H,
\]
where \(C_H\) is the coordinate-opening handle for the Boolean tensor \(H\).
The final translation in Step A30 replaces \(H\) by the lane index
\(q_{\RAA}(H)\), but the logical requests are easier to read before that
packing step.

\begin{description}[leftmargin=0.25in,itemsep=0.08em,topsep=0.1em,parsep=0pt]
  \item[Value request.]
    \((m_\star,C_{m_\star},z_{\mathrm{in}},v_m)\), with
    \(m_\star\in\F^{\B^k}\).  The verifier also checks the scalar equality
    \(v_m=v_\star\).
  \item[Accumulator endpoint requests.]
    \((u_3,C_{u_3},r_{\mathrm{chk}},s_1)\) and
    \((c_\star,C_{c_\star},r_{\mathrm{chk}},s_2)\), with
    \(u_3,c_\star\in\F^{\B^n}\).  These are the output-side MLE claims of the
    two accumulator reductions.
  \item[Permutation root requests.]
    The first permutation contributes
    \[
      (f_1,C_{f_1},\mathsf{root},v_{f_1,\mathrm{root}}),
      \qquad
      (g_1,C_{g_1},\mathsf{root},v_{g_1,\mathrm{root}}).
    \]
    The second permutation contributes the analogous requests for \(f_2\) and
    \(g_2\).  The tensors \(f_i,g_i\) live in \(\F^{\B^{n+1}}\).  The verifier
    checks the two scalar root equalities
    \(v_{f_i,\mathrm{root}}=v_{g_i,\mathrm{root}}\).
  \item[Permutation leaf requests.]
    For the first permutation:
    \[
      (u_1,C_{u_1},r_{\mathrm{chk}},v_{u_1}),\quad
      (u_2,C_{u_2},r_{\mathrm{chk}},v_{u_2}),\quad
      (f_1,C_{f_1},\mathsf{leaf}(r_{\mathrm{chk}}),v_{f_1}),\quad
      (g_1,C_{g_1},\mathsf{leaf}(r_{\mathrm{chk}}),v_{g_1}).
    \]
    The second permutation contributes the analogous four requests with
    \((u_3,u_4,f_2,g_2,\widehat\pi_2)\).  The verifier checks the scalar leaf
    equations from Step CP4.
  \item[Six-sumcheck terminal requests.]
    The terminal point \(r_{\mathrm{sc}}\in\F^n\) contributes
    \((u_2,C_{u_2},r_{\mathrm{sc}},v_1)\) and
    \((u_4,C_{u_4},r_{\mathrm{sc}},v_2)\).  For each
    \(T\in\{f_1,g_1,f_2,g_2\}\), it also contributes
    \[
      (T,C_T,\mathsf{parent}(r_{\mathrm{sc}}),v_{T,p}),\quad
      (T,C_T,\mathsf{left}(r_{\mathrm{sc}}),v_{T,l}),\quad
      (T,C_T,\mathsf{right}(r_{\mathrm{sc}}),v_{T,r}).
    \]
\end{description}

The public factors \(\widehat A(r_{\mathrm{chk}},r_{\mathrm{sc}})\),
\(\chi_{r_{\mathrm{chk}}}(r_{\mathrm{sc}})\),
\(\widehat\pi_i(r_{\mathrm{chk}})\), \(\Idx(r_{\mathrm{chk}})\), and the powers
of the batching scalar \(\lambda\) are not authentication requests.  They are
known field elements used in scalar checks after the requested MLE values have
been authenticated.

For the first permutation, let \(v_{f_1}\) denote the claimed value of
\(\MLE(f_1)(\mathsf{leaf}(r_{\mathrm{chk}}))\).  The virtual \(u_1\) handle
means that the leaf equation is concretely
\[
  v_{f_1}
  =
  \alpha-\bigl(\MLE(m_\star)(\mathsf{msg}(r_{\mathrm{chk}}))
  +\beta\,\widehat\pi_1(r_{\mathrm{chk}})\bigr).
\]
The value \(\MLE(m_\star)(\mathsf{msg}(r_{\mathrm{chk}}))\) is not read by the
verifier.  It is one of the MLE claims authenticated through
\(\MLEAuth\), using the virtual coordinate handle \(C_{u_1}\).

\subsection{Top-Level Protocol}

The PCS has one commitment phase and one opening phase.  The commitment phase
binds the encoded matrix.  The opening phase first checks the original claimed
evaluation using slice values, then folds those slice claims into one RAA claim,
and finally invokes the compiled RAA adapter to authenticate the resulting MLE
evaluation requests.

% DO NOT MODIFY: Frozen semantic anchor.
% This top-level PCS protocol is the current source of truth for Blaze protocol semantics.
\begin{protocolbox}{Protocol: Blaze PCS}{fig:protocol-blaze-pcs}
\scriptsize\raggedright
\noindent\textbf{DO NOT MODIFY.} Frozen semantic anchor; this figure is the
current source of truth for Blaze protocol semantics.

\smallskip
\noindent Public parameters: \(\ell,k,n\), \(L=2^\ell\), \(K=2^k\),
\(N=2^n\), the RAA map
\(\RAA:\F^{\B^k}\to\F^{\B^n}\), \(\mathsf{RowFold}\) from
Figure~\ref{fig:protocol-row-fold-handle}, the compiled RAA check from
Figure~\ref{fig:protocol-compiled-raa-check}, and
\(\mathcal F_{\mathrm{MsgOracle}}\).

\smallskip
\noindent\(\underline{\mathsf{Blaze.Commit}(\verifier:\bot,\ \prover:m\in\F^{\B^\ell\times\B^k})\to C_m\in\Handle}\)

\smallskip
\begin{enumerate}[leftmargin=0.24in,label=\textbf{C\arabic*.},itemsep=0.06em,topsep=0.06em,parsep=0pt]
  \item \(\prover\): for each \(h\in\B^\ell\), compute
        \(c[h,\ast]\gets\RAA(m[h,\ast])\in\F^{\B^n}\).
  \item \(\prover\): define the column tensor \(C\in(\F^{\B^\ell})^{\B^n}\)
        by \(C[b]=c[\ast,b]\) for \(b\in\B^n\).
  \item \(\prover,\verifier\): \(R_c\gets
        \mathcal F_{\mathrm{MsgOracle}}.\mathsf{Submit}(\prover:C,\verifier:\bot)\).
  \item \(\mathsf{Blaze}\): return \(C_m=R_c\).
\end{enumerate}

\smallskip
\noindent\(\underline{\mathsf{Blaze.Open}(\verifier:(C_m,z=(z_{\mathrm{slice}},z_{\mathrm{in}}),v),\ \prover:(m,C_m,z,v))\to(\verifier:a\in\{0,1\})}\)

\smallskip
\begin{enumerate}[leftmargin=0.24in,label=\textbf{O\arabic*.},itemsep=0.06em,topsep=0.06em,parsep=0pt]
  \item \(\prover\): for each \(h\in\B^\ell\), compute
        \(u[h]\gets\MLE(m[h,\ast])(z_{\mathrm{in}})\).
  \item \(\prover\to\verifier\): \(u\in\F^{\B^\ell}\).
  \item \(\verifier\): if
        \(v\ne\sum_{h\in\B^\ell}\chi_h(z_{\mathrm{slice}})u[h]\), return
        \((\verifier:0)\).
  \item \(\verifier\): sample \(\rho\leftarrow\F^{\B^\ell}\).
  \item \(\verifier\to\prover\): \(\rho\).
  \item \(\prover,\verifier\): \(R_\star\gets
        \mathsf{RowFold.Commit}(\verifier:(C_m,\rho),\prover:(C_m,\rho))\).
  \item \(\prover\): compute \(m_\star\in\F^{\B^k}\) by
        \(m_\star[a]\gets\sum_{h\in\B^\ell}\rho[h]m[h,a]\).
  \item \(\verifier\): compute
        \(v_\star\gets\sum_{h\in\B^\ell}\rho[h]u[h]\).
  \item \(\prover\): choose \(c_\star^{\mathrm{wit}}\in\F^{\B^n}\).
  \item \(\prover,\verifier\): \(a_{\RAA}\gets
        \mathsf{CompiledRAA.Check}
        (\verifier:(R_\star,z_{\mathrm{in}},v_\star),
        \prover:(R_\star,m_\star,c_\star^{\mathrm{wit}}))\).
  \item \(\verifier\): return \((\verifier:a_{\RAA})\).
\end{enumerate}
\end{protocolbox}

The tensor \(h_{c_\star}=c_\star^{\mathrm{wit}}\) is only the prover's witness
candidate for the virtual folded word.  Its authenticated coordinate handle is
\(R_\star\), so backend openings of \(c_\star\) are checked against the
previously committed matrix \(c\) through Figure~\ref{fig:protocol-row-fold-handle}.
There is not a second commitment to \(c_\star^{\mathrm{wit}}\).  The witness
tensor supplies prover-side values for algebraic computation; the verifier-side
coordinate source for \(c_\star\) is the virtual row-fold handle \(R_\star\).

\subsection{What Soundness Uses}

The top-level slice check says that the original claim follows from the slice
values \(u[h]\).  The random vector \(\rho\) then turns many slice claims into
one folded claim
\[
  \MLE(m_\star)(z_{\mathrm{in}})=v_\star.
\]
At the same time, the row-fold handle turns the committed encoded matrix into
one openable virtual tensor \(c_\star\).  If all committed slices were honest,
then \(c_\star=\RAA(m_\star)\).  If not, the random fold makes the disagreement
survive in the folded object except with the usual small probability over
\(\rho\).  This is the interleaving trick in its most concrete form: a query to
the folded word \(c_\star[b]\) is emulated by opening the committed column
\(c[\ast,b]\) and taking the verifier's random linear combination.
Before the backend proximity theorem is applied, this statement is only about
the verifier's openable virtual word.  It should not be read as saying that the
virtual word is already known to be a valid RAA codeword.

The native RAA compiler does not prove proximity.  It produces algebraic
evaluation requests whose truth would imply the RAA stage relations for the
folded statement.  The \(\MLEAuth\) layer authenticates those requests against
the coordinate-opening handles, using the BaseFold/RMLE backend described in
Figure~\ref{fig:protocol-mle-auth-basefold}.  Few coordinate openings do not by
themselves prove that an opened word is a valid encoding; the backend proximity
argument and the code distance are what make the authenticated MLE claims mean
one nearby tensor rather than unrelated sampled answers.

In an accepting transcript, the backend extractor should be read as extracting
a nearby valid backend object for each authenticated handle.  For the folded
message and codeword handles, write the extracted tensors as
\(m_\star^{\mathrm{ext}}\) and \(c_\star^{\mathrm{ext}}\).  The algebraic RAA
checks prove the relation for the extracted objects:
\[
  c_\star^{\mathrm{ext}}=\RAA(m_\star^{\mathrm{ext}}),
  \qquad
  \MLE(m_\star^{\mathrm{ext}})(z_{\mathrm{in}})=v_\star.
\]
This does not say that the verifier-facing virtual tensor \(c_\star\) is
literally equal to \(\RAA(m_\star^{\mathrm{ext}})\) at every coordinate.  It says that
\(c_\star\) is bound by \(R_\star\), authenticated at queried coordinates, and
forced close enough to the extracted code-consistent tensor for the backend
soundness theorem to apply.

The soundness case split is therefore explicit.  First, if the packed tensor
submitted to \(\MLEAuth\) is far from the backend language, the BaseFold/RMLE
proximity check is supposed to reject except with its proximity error.  Second,
if the packed tensor is close enough to an extracted tensor but some request
\((q_i,r_i,v_i)\in\mathcal E_{\RAA}\) is false for that extracted tensor, then
the batching challenge in Figure~\ref{fig:protocol-mle-eval-batch} must either
accidentally cancel the false requests or produce a false Boolean-sum claim.
Third, if the batched Boolean-sum claim is false, sumcheck must be fooled.  At
the terminal point, the final backend proof authenticates the one remaining
claim
\[
  \MLE(A)(r_{\mathrm{batch}})=v_{\mathrm{batch}},
\]
except in the endpoint-degenerate case where the public weight
\(W_\gamma(r_{\mathrm{batch}})\) vanishes and the terminal equation no longer
uses that value.  That degeneracy is part of the batching soundness accounting,
not a new RAA condition.  The Blaze algebra supplies the requests; the backend
supplies the meaning of those requests against one bound nearby tensor.

Thus the final proof is a composition:
\[
  \begin{aligned}
  &\text{slice evaluation}
  +\text{ random row fold}\\
  &\quad+\text{ native RAA algebra}
  +\text{ authenticated MLE requests}.
  \end{aligned}
\]
